\section{Интеграл Лебега}
\subsection{Интеграл Лебега для неотрицательных функций}
\paragraph{Лекция 9.}

Пусть даны $X = (X, \mathfrak{M}, \mu)$.
Напомним
\begin{definition}
    $h: X \to \R$ называется \textit{простой}, если $h(x) = \sum\limits_{k=1}^{N} a_k \chi_{E_k}(x)$, где $a_1, \dots, a_N \in \R$,
    $E_1, \dots, E_N \in \mathfrak{M}$.
    
    Также $h$ — простая функция тогда и только тогда, когда
    \[X = \bigsqcup\limits_{k=1}^{N'} G_k, \ h(x) = \sum\limits_{k=1}^{N'}\alpha_k \chi_{G_k}(x).\]
\end{definition} 
\begin{definition}
    Пусть $h: X \to [0, +\infty)$ — простая функция.
    Пусть $E \in \mathfrak{M}$.
    \textit{Интеграл Лебега от $h$ по множеству $E$}:
    \[\int_E h(x) d\mu(x) := \sum\limits_{k=1}^{N'}\alpha_k \mu(E \cap G_k).\]
\end{definition} 

\begin{lemma}
    Определение корректно, то есть не зависит от представления $h$.
\end{lemma} 
\begin{proof}
    Пусть имеем два разбиения $X$:
    \[X = \bigsqcup\limits_{i=1}^{N_1} A_i, \tab X = \bigsqcup\limits_{j=1}^{N_2}B_j.\]
    \[h(x) = \sum\limits_{i=1}^{N_1}\chi_{A_i} (x) \cdot a_i,
    \tab
    h(x) = \sum\limits_{j=1}^{N_2} \chi_{B_j}(x) \cdot b_j.\]
    Определим более мелкое разбиение
    \[C_{ij} := A_i \cap B_j,\]
    тогда $C_{ij} \neq C_{i'j'}$ при $(i,j) \neq (i', j')$.
    Тогда
    \[X = \bigsqcup\limits_{i=1}^{N_1}\bigsqcup\limits_{j=1}^{N_2}C_{ij}.\]

    Пусть $E \in \mathfrak{M}$.
    \begin{multline*}
        \int_E h(x) d\mu(x) 
        = \sum\limits_{i=1}^{N_1}a_i \mu(A_i \cap E) 
        = \left| A_i \cap E 
        = \bigsqcup\limits_{j=1}^{N_1} C_{ij} \cap E \right| 
        =\\= 
        \sum\limits_{i=1}^{N_1} \sum\limits_{j=1}^{N_2} a_i \mu(E \cap C_{ij})
        = \left| h(x) = a_i = b_j \text{ при } x \in A_i \cap B_j \right| 
        =\\=
        \sum\limits_{i=1}^{N_1} \sum\limits_{j=1}^{N_2} b_j \mu(E \cap C_{ij}) 
        = \sum\limits_{j=1}^{N_2}b_j \underbrace{\sum\limits_{i=1}^{N_1} \mu(E \cap C_{ij})}_{=\mu(E \cap B_j)} 
        =\\=
        \sum\limits_{j=1}^{N_2} b_j \mu(E \cap B_j) = \int_E h(x) d\mu(x).
    \end{multline*} 
\end{proof} 

\begin{definition}
    Пусть теперь $f: X \to [0, +\infty]$ — измеримая на $X$ функция.
    Тогда
    \[\int_E f(x) d\mu(x) := \sup \left\{\int_E h(x) d\mu(x) \ | \ h \text{ — простая, } 0 \leqslant h(x) \leqslant f(x) \ \forall x \in E\right\}.\]
\end{definition} 
\begin{remark}
    На простых функциях это определение переходит в предыдущее.
\end{remark}
\begin{statement}
    Пусть $\widetilde{h}(x)$ и $h(x)$ — простые функции и
    $\widetilde{h}(x) \leqslant h(x)$ на $E$.
    Тогда $$\int_E \widetilde{h}(x) d\mu(x) \leqslant \int_E h(x) d\mu(x).$$
\end{statement} 
\begin{proof}
    Пусть $X = \bigsqcup\limits_{k=1}^{N}C_k$ — общее разбиение и для $h$, и для $\widetilde{h}$.
    Тогда при $0 \leqslant \widetilde{c}_k \leqslant c_k$:
    \[h(x) = \sum\limits_{k=1}^{N} c_k \chi_{C_k}(x), \tab \widetilde{h}(x) = \sum\limits_{k=1}^{N} \widetilde{c}_k \chi_{C_k}(x).\]
    В силу предыдущей леммы
    \[\int_E \widetilde{h}(x) d\mu(x) 
    = \sum\limits_{k=1}^{N}\widetilde{c}_k \mu(A_k \cap E)
    \leqslant \sum\limits_{k=1}^{N} c_k \mu(A_k \cap E)
    = \int_E h(x) d\mu(x)\]
    %\[\implies \sup\limits_{0 \leqslant \widetilde{h} \leqslant h}
    %\int_E \widetilde{h}(x) d\mu(x) \leqslant \int_E h(x) d\mu(x).%\]
\end{proof} 

\begin{definition}
    Пусть $f: X \to \R$ — измеримая функция.
    Пусть $E \in \mathfrak{M}$.
    Определим
    \[f_+(x) := \max\{0, f(x)\} \geqslant 0,\]
    \[f_-(x) := \max\{0, -f(x)\} \geqslant 0.\]
    Тогда если 
    \[\begin{cases}
        \int_E f_+(x) d\mu(x) < +\infty, \tab (1)\\
        \int_E f_-(x) d\mu(x) < +\infty, \tab (2)
    \end{cases}
    \]
    то \textit{$f$ интегрируема по Лебегу на множестве $E$}, а сам интеграл обозначается
    \[\int_E f(x) d\mu(x) := \int_E f_+(x) d\mu(x) - \int_E f_-(x) d\mu(x).\]
\end{definition} 
\begin{remark}
    Если хотя бы одно из условий (1)-(2) не выполнено, то $f$ называется неинтегрируемой по Лебегу на $E$.
\end{remark}

\subsection{Свойства интеграла Лебега для неотрицательных функций}
\begin{statement}[Монотонность по функциям]
    Пусть $f,g: X \to [0, +\infty]$ — измеримые функции, $E \in \mathfrak{M}$.
    Пусть $f(x) \leqslant g(x)$ для любого $X \in E$.
    Тогда
    \[\int_E f(x) d\mu(x) \leqslant \int_E g(x) d\mu(x).\]
\end{statement} 
\begin{proof}
    Для простых функций доказано в утверждении выше.
    Пусть $h$ — простая неотрицательная функция:
    \begin{equation}
        0 \leqslant h(x) \leqslant f(x) \leqslant g(x) \ \forall x \in E. \label{eq:9.1}
    \end{equation}
    Тогда по определению интеграла Лебега
    \[\int_E h(x) d\mu(x) \leqslant \int_E g(x) d\mu(x).\]
    Возьмём супремум от обеих частей по $h$, удовлетворяющим \eqref{eq:9.1}, тогда
    \[\int_E f(x) d\mu(x) \leqslant \int_E g(x) d\mu(x).\]
\end{proof} 

\begin{statement}[Умножение на характеристическую функцию]
    Пусть $f: X \to [0, +\infty]$ — измеримая функция,
    $E \in \mathfrak{M}$ — измеримое множество.
    Тогда
    \[\int_E f(x) d\mu(x) = \int_X \chi_E(x) f(x) d\mu(x).\]
\end{statement} 
\begin{proof}
    Сначала докажем для простых функций.
    $$h(x) := \sum\limits_{k=1}^{N} a_k \chi_{A_k}(x),$$
    где $A_1, \dots, A_N$ — разбиение $X$.
    \[\chi_E h(x) := \sum\limits_{k=1}^{N} a_k \chi_{A_k \cap E}(x) + 0 \cdot \chi_{X \setminus E}(x).\]
    Тогда
    \[A_1 \cap E, \dots, A_N \cap E, X \setminus E\]
    — разбиение $X$ для функции $\chi_E h$.
    \begin{multline*}
        \int_X (h \chi_E)(x) d\mu(x) 
        = a_1 \cdot \mu(A_1 \cap E) + \dots + a_N \cdot \mu(A_N \cap E) + 0 \cdot \mu(X \setminus E)
        =\\=
        \sum\limits_{i=1}^{N} a_i \mu(A_i \cap E)
        = \int_E h(x) d\mu(x).
    \end{multline*} 
    Пусть $f: X \to [0, +\infty]$ — неотрицательная измеримая функция, которая не является простой.
    \[0 \leqslant h(x) \leqslant (\chi_E f)(x) \ \forall x \in X
    \text{ и }
    h(x) = (\chi_E h)(x).\]
    В силу только что доказанного для простых функций:
    \[\int_X h(x) d\mu(x) 
    = \int_X \chi_E(x) h(x) d\mu(x)
    = \int_E h(x) d\mu(x)
    \leqslant \int_E f(x) d\mu(x).\]
    Взяв супремум по всем $h$, удовлетворяющим \eqref{eq:9.1}, получим
    \[\int_X \chi_E(x) f(x) d\mu(x) \leqslant \int_E f(x) d\mu(x).\]
    Пусть $g$ — простая неотрицательная функция, определённая на всём $X$.
    Пусть
    \begin{equation}
        0 \leqslant g(x) \leqslant f(x) \quad \forall x \in E \label{eq:9.2}
    \end{equation}
    Поскольку $g$ — простая, то в силу доказанного выше
    \[\int_E g(x) d\mu(x) 
    = \int_X (\chi_E g)(x) d\mu(x)
    \leqslant \int_X (\chi_E f)(x) d\mu(x).\]
    Взяв супремум от обеих частей по $g$, удовлетворяющим \eqref{eq:9.2}, получим
    \[\int_E f(x) d\mu(x) \leqslant \int_X (\chi_E f) (x) d\mu(x).\]
\end{proof} 

\begin{statement}[Монотонность по множеству]
    $A,B \in \mathfrak{M}$ и $A \subset B$,
    $f: X \to [0, +\infty]$ — измеримая функция.
    Тогда
    \[\int_A f(x) d\mu(x) \leqslant \int_B f(x) d\mu(x).\]
\end{statement} 
\begin{proof}
    \[A \subset B
    \implies
    \chi_A f(x) \leqslant \chi_B f(x) \quad \forall x \in X,\]
    тогда в силу монотонности по функции
    \[\underbrace{\int_X (\chi_A f)(x) d\mu(x)}_{\int_A f(x) d\mu(x)} 
    \leqslant 
    \underbrace{\int_X (\chi_B f)(x) d\mu(x)}_{\int_B f(x) d\mu(x)}.\]
    Следовательно,
    \[\int_A f(x) d\mu(x) \leqslant \int_B f(x) d\mu(x).\]
\end{proof} 

\begin{statement}[Неравенство Чебышёва]
    Пусть $f$ — интегрируемая на $X$ и неотрицательная на нём функция.
    Тогда для любого $t>0$:
    \[E_t := \{x \in X \ | \ f(x) \geqslant t\},\]
    \[\mu(E_t) \leqslant \frac{1}{t} \int_X f(x) d\mu(x).\]
\end{statement} 
\begin{proof}
    \[t \mu(E_t)
    = \int_{E_t} t \chi_E(x) d\mu(x)
    \leqslant \int_{E_t} f(x) d\mu(x)
    \leqslant \int_X f(x) d\mu(x).\]
    Поделив обе части на $t$, получим утверждение.
\end{proof} 

\begin{statement}
    Если $f: X \to [0, +\infty]$ — интегрируемая по Лебегу функция, то $f(x) < +\infty$ для $\mu$-почти всех $x \in X$.
\end{statement} 
\begin{proof}
    $$E_{\infty} := \{x \in X \ | \ f(x) = +\infty\}.$$
    \[E_{\infty} = \bigcap\limits_{t>0} E_t.\]
    По неравенству Чебышёва
    \[\mu(E_t) \leqslant \frac{1}{t} \underbrace{\int_X f(x) d\mu(x)}_{=C},\]
    тогда в силу монотонности меры
    \[0 \leqslant \mu(E_\infty) \leqslant \frac{C}{t} \ \forall t \geqslant 0
    \implies
    \mu(E_\infty) = 0.\]
\end{proof} 

\begin{statement}
    Пусть $f: X \to [0, +\infty]$ — измеримая функция, $\mu(E) > 0$, $E \in \mathfrak{M}$, $f(x) > 0$ для любого $x \in E$.
    Тогда 
    \[\int_E f(x) d\mu(x) > 0.\]
\end{statement} 
\begin{proof}
    \[\exists n_0 \in \N: \ \mu\left(\underbrace{\left\{x \in E \ | \ f(x) \geqslant \frac{1}{n_0}\right\}}_{E_0}\right) > 0.\]
    Для $\chi_{E_0} f$ применим неравенство Чебышёва:
    \[0 < \frac{\mu(E_0)}{n_0} 
    \leqslant \int_X (\chi_{E_0}f)(x) d\mu(x)
    \leqslant \int_X (\chi_Ef)(x) d\mu(x)
    = \int_E f(x) d\mu(x).\]
\end{proof} 

\begin{statement}
    Пусть $\mu(E) = 0$, $E \in \mathfrak{M}$, $f: E \to [0, +\infty]$ — измеримая неотрицательная функция.
    Тогда
    \[\int_E f(x) d\mu(x) = 0.\]
\end{statement} 
\begin{proof}
    Для простых функций утверждение очевидно.
    В общем случае супремум от нуля есть ноль.
\end{proof} 

\begin{theorem}[Леви\footnote{Беппо Леви (1875-1961) — итальянский математик.}]
    Пусть $\{f_n\}$ — последовательность неотрицательных измеримых на $X$ функций:
    \[0 \leqslant f_1(x) \leqslant \dots \leqslant f_n(x) \leqslant \dots.\]
    Тогда 
    \[\lim\limits_{n\to \infty}\int_X f_n(x) d\mu(x) = \int_X \lim\limits_{n \to \infty}f_n(x) d\mu(x),\]
    \[\exists f(x) = \lim\limits_{n\to \infty}f_n(x) \in [0, +\infty] \ \forall x \in X, \ f \text{ измерима на } X.\]
\end{theorem} 
\begin{proof}
    Поскольку $\{f_n\}$ — монотонная последовательность, то при
    \[L_n := \int_X f_n(x) d\mu(x), \ n \in \N.\]
    последовательность $\{L_n\} \subset [0, +\infty]$ тоже монотонна, следовательно,
    \[\exists \lim\limits_{n\to \infty}L_n := L \in [0, +\infty].\]
    В силу монотонности интеграла по функции
    \[\int_X f_n(x) d\mu(x) \leqslant \int_X f(x) d\mu(x) \ \forall n \in \N
    \implies
    L \leqslant \int_X f(x) d\mu(x).\]
    Зафиксируем произвольное число $\theta \in (0, 1)$.
    Определим простую функцию $g$: $0 \leqslant g(x) \leqslant f(x)$, тогда $\theta g(x) < f(x)$.
    \[X_n := X(f_n > \theta g), \ n \in \N \text{ и } X_{n+1} \supset X_n \ \forall n \in \N.\]
    Кроме того, $X = \bigcup\limits_{n=1}^{\infty} X_n$.
    Если $g(x) = 0$, то $x \in \bigcup\limits_{n=1}^{\infty} X_n$, поскольку $f_n(x) > \theta g(x)$ для любого $x \in X_n$.
    Если же $\theta g(x) > 0$, то по определению предела поскольку \[\lim\limits_{n \to \infty} f_n(f) \geqslant g(x) > \theta g(x),\]
    то существует 
    $$N(x) \in \N: f_n(x) > \theta g(x) \ \forall n \geqslant N(x)
    \implies
    x \in \bigcup\limits_{n \geqslant N(x)}X_n.$$
    В силу непрерывности меры снизу если $A \in \mathfrak{M}$, то $\mu(A) = \lim\limits_{n\to \infty}\mu(A \cap X_n)$.

    Поскольку $g$ — простая функция, то есть $X = \bigsqcup\limits_{k=1}^{N}A_k$ — допустимое разбиение $X$.
    \[g(x) = \sum\limits_{k=1}^{N}a_k \chi_{A_k}(x) \ \forall x \in X.\]
    \begin{multline*}
        \int_X \theta g(x) d\mu(x) 
        = \sum\limits_{k=1}^{N} \theta a_k \mu(A_k)
        = \sum\limits_{k=1}^{N} \lim\limits_{n\to \infty} \theta a_k \mu(A_k \cap X_n)
        =\\= 
        \lim\limits_{n \to \infty} \sum\limits_{k=1}^{N} \theta a_k \mu(A_k \cap X_n) 
        = \lim\limits_{n\to \infty} \int_{X_n} \theta g(x) d\mu(x)
        \leqslant\\\leqslant 
        \lim\limits_{n\to \infty}\int_{X_n} f_n(x) d\mu(x)
        \leqslant \lim\limits_{n\to \infty} \int_X f_n(x) d\mu(x) = L.
    \end{multline*} 
    Получили $\int_X \theta g(x) d\mu(x) \leqslant L$ — это верно для любого $\theta \in (0,1)$ и для любого $g$: $0 \leqslant g \leqslant f$.
    Сначала возьмём супремум по $\theta$:
    \[\int_X g(x) d\mu(x) \leqslant L,\]
    а потом берём супремум по $g$.
\end{proof} 

\begin{lemma}[Фату]
    Пусть $\{f_n\}$ — последовательность неотрицательных измеримых на $X$ функций.
    Тогда
    \[\int_X \lowlim\limits_{n\to \infty} f_n(x) d\mu(x)
    \leqslant
    \lowlim\limits_{n\to \infty} \int_X f_n(x) d\mu(x).\]
\end{lemma} 
\begin{proof}
    \[\lowlim\limits_{n\to \infty} f_n(x) = \sup\limits_{n \in \N} \underbrace{\inf\limits_{k\geqslant n} f_k(x)}_{g_n(x)} = \lim\limits_{n\to \infty} g_n(x) \ \forall x \in X\]
    Тогда $\{g_n\}$ — последовательность таких измеримых неотрицательных функций, что $g_{n+1}(x) \geqslant g_n(x)$ для любых $x \in X$.
    
    Применим к $\{g_n\}$ теорему Леви:
    \[\underbrace{\int_X \lim\limits_{n\to \infty} g_n(x) d\mu(x)}_{\int_X \lowlim\limits_{n\to \infty} f_n(x) d\mu(x)} 
    =
    \underbrace{\lim\limits_{n\to \infty} \int_X g_n(x) d\mu(x)}_{\lim\limits_{n\to \infty} \int_X \inf\limits_{k \geqslant n} f_k(x) d\mu(x)},\]
    а поскольку
    \[\int_X \inf\limits_{k \geqslant n} f_k(x) d\mu(x) \leqslant \inf\limits_{k \geqslant n} \int_X f_k(x) d\mu(x),\]
    то
    \begin{multline*}
        \int_X \lowlim\limits_{n\to \infty} f_n(x) d\mu(x)
        =
        \lim\limits_{n\to \infty} \int_X \inf\limits_{k \geqslant n} f_k(x) d\mu(x)
        \leqslant\\\leqslant
        \lim\limits_{n\to \infty}\inf\limits_{k \geqslant n} \int_X f_k(x) d\mu(x)
        \leqslant
        \lowlim\limits_{n\to \infty} \int_X f_n(x) d\mu(x).
    \end{multline*} 
\end{proof} 



\paragraph{Лекция 10.}
\begin{lemma}[Аддитивность]
    Пусть $f,g$ — неотрицательные измеримые функции.
    Тогда
    \[\int_X f(x) + g(x) d\mu(x)
    =
    \int_X f(x) d\mu(x) + \int_X g(x)d\mu(x).\]
\end{lemma} 
\begin{proof}
    Пусть $\psi$, $h$ — такие простые функции, что:
    \[0 \leqslant \psi \leqslant f,\]
    \[0 \leqslant h \leqslant g.\]
    \[\psi(x) = \sum\limits_{k=1}^{N_1} \widetilde{a}_k \chi_{E_k}(x),\]
    \[h(x) = \sum\limits_{k=1}^{N_2} \widetilde{b}_k \chi_{G_k}(x).\]
    Пусть $\{C_k\}_{k=1}^N$ — общее разбиение для $\psi$ и $h$, то есть 
    \[X = \bigsqcup\limits_{k=1}^{N}C_k\]
    и
    \[\psi(x) = \sum\limits_{k=1}^{N}a_k \chi_{C_k}(x),\]
    \[h(x) = \sum\limits_{k=1}^{N}b_k \chi_{C_k}(x).\]
    Тогда
    \begin{multline}
        \int_X \psi(x) + h(x) d\mu(x)
        = \sum\limits_{k=1}^{N}(a_k+b_k) \mu(C_k)
        =\\= \sum\limits_{k=1}^{N}a_k\mu(C_k) + \sum\limits_{k=1}^{N}b_k \mu(C_k)
        = \int_X \psi(x) d\mu(x) + \int_X h(x) d\mu(x). \label{eq:9.3}
    \end{multline} 
    В общем случае строим последовательность неотрицательных простых функций, которые всюду на $X$ сходятся к $f$ и $g$:
    \[\psi_n \xrightarrow{X} f, \ n \to \infty,\]
    \[h_n \xrightarrow{X} g, \ n \to \infty.\]
    Более того, можно считать $\{\psi_n\}$ и $\{h_n\}$ монотонны.
    Тогда запишем \eqref{eq:9.3} для $\psi_n$ и $h_n$ и, пользуясь теоремой Беппо Леви, перейдём к пределу при $n\to \infty$.
\end{proof} 

\begin{lemma}[Положительная однородность]
    Пусть $f \geqslant 0$ — измеримая на $X$ функция.
    Тогда для любого $\alpha \geqslant 0$
    \[\int_X \alpha \cdot f d\mu(x) = \alpha \cdot \int_X f d\mu(x).\]
\end{lemma} 
\begin{proof}
    Доказательство аналогично предыдущей лемме: сначала доказываем для простых функций, а потом, пользуясь теоремой Беппо Леви, переходим к пределу.
\end{proof} 

\begin{lemma}
    Пусть $E \in \mathfrak{M}$ и $\{E_k\}_{k=1}^N \subset \mathfrak{M}$ и $E = \bigcup\limits_{k=1}^{N}E_k$.
    Тогда для произвольной измеримой и неотрицательной функции $f$ верно следующее:
    \[\int_E f d\mu(x) < +\infty
    \Longleftrightarrow
    \int_{E_k} f d\mu(x) < +\infty \tab \forall k \in \{1, \dots, N\}.\]
\end{lemma} 
\begin{proof}\tab
    \begin{itemize}
        \item[\underline{$\implies$}]
        \[\int_{E_k} f d\mu(x) \leqslant \int_E f d\mu(x) < +\infty \tab \forall k \in \{1, \dots, N\}.\]
        \item[\underline{$\impliedby$}]
        Пусть
        \[\int_{E_k} f d\mu(x) < +\infty \tab \forall k \in \{1, \dots, N\}.\]
        Тогда
        \begin{multline*}
            \int_E f d\mu(x) = \int_X f \chi_E d\mu(x)
            \leqslant
            \int_X \sum\limits_{k=1}^{N} f \chi_{E_k} d\mu(x)
            =\\= \sum\limits_{k=1}^{N} \int_X f \chi_{E_k} d\mu(x)
            = \sum\limits_{k=1}^{N} \int_{E_k} f d\mu(x) < +\infty.
        \end{multline*} 
    \end{itemize}
\end{proof} 

\begin{remark}
    В условиях предыдущей леммы если
    \[E = \bigsqcup\limits_{k=1}^{N}E_k,\]
    то
    \[\int_E f d\mu(x) = \sum\limits_{k=1}^{N} \int_{E_k} f d\mu(x),\]
    поскольку
    \[\chi_E = \sum\limits_{k=1}^{N} \chi_{E_k}\]
    и можно воспользоваться линейностью интеграла по функции.
\end{remark}

\begin{lemma}
    Пусть $E \in \mathfrak{M}$ и $\mu(E) > 0$.
    Пусть $f$ — измеримая неотрицательная функция и $f>0$ на $E$.
    Тогда
    \[\int_E f d\mu > 0.\]
\end{lemma} 
\begin{proof}
    Для любого $n \in \N$
    \[E_n := E\left(f > \frac{1}{n}\right) \text{ и } E = \bigcup\limits_{n=1}^{\infty} E_n \implies\]
    \[\implies \exists n_0 \in \N: \ \mu(E_{n_0}) > 0 \implies\]
    \[\implies \int_E f d\mu \geqslant \int_{E_{n_{0}}} f d\mu
    \geqslant \int_{E_{n_0}} \frac{1}{n_0} d\mu = \frac{\mu(E_{n_0})}{n_0} > 0. \]
\end{proof} 


\subsection{Свойства интеграла Лебега от произвольных измеримых функций}
\begin{statement}
    $f$ интегрируема по Лебегу тогда и только тогда, когда $|f|$ интегрируема.
\end{statement} 
\begin{proof}\tab
    \begin{itemize}
        \item[\underline{$\implies$}]
        \[f_+ := \max\{f,0\},\]
        \[f_- := \max\{-f,0\}.\]
        По определению $f$ интегрируема по Лебегу тогда и только тогда, когда
        \[\begin{cases}
            \int_X f_+(x) d\mu(x) < +\infty, \tab (1)\\
            \int_X f_-(x) d\mu(x) < +\infty, \tab (2)
        \end{cases}
        \]
        Поскольку
        \[|f| = f_+ + f_-,\]
        то можно сказать, что если выполнены (1)-(2), то $|f|$ интегрируема.

        \item[\underline{$\impliedby$}]
        Пусть $|f|$ интегрируема.
        Тогда
        \[
        \begin{cases}
            f_+ \leqslant |f|,\\
            f_- \leqslant |f|,
        \end{cases}
        \]
        следовательно, по свойствам интеграла Лебега для неотрицательных функций
        $f_+$ и $f_-$ интегрируемы по Лебегу, поэтому интегрируема по Лебегу и $f$.
    \end{itemize}
\end{proof} 

\begin{statement}
    Пусть $f$ интегрируема по Лебегу на $E$.
    Тогда
    \[\left|\int_E f d\mu\right| \leqslant \int_E |f| d\mu.\]
\end{statement} 
\begin{proof}
    \begin{multline*}
        \left|\int_E f d\mu\right|
        = \left|\int_E f_+ d\mu - \int_E f_- d\mu\right|
        \leqslant\\\leqslant
        \left|\int_E f_+ d\mu\right| + \left|\int_E f_- d\mu\right|
        =\\= \int_E f_+ d\mu + \int_E f_- d\mu
        = \int_E |f| d\mu.
    \end{multline*}
\end{proof} 

\begin{statement}
    Пусть $f$ интегрируема по Лебегу на $E$.
    Тогда $|f(x)| < +\infty$ для $\mu$-почти всех $x \in E$.
\end{statement} 
\begin{proof}
    Поскольку $f$ интегрируема по Лебегу тогда и только тогда, когда $|f|$ интегрируема по Лебегу, то по условию
    \[\int_E |f| d\mu < +\infty.\]
    Определим
    \[E_0 := \{x \in E \ | \ f(x) = \pm \infty\}.\]
    \[+\infty > \int_E |f|d\mu \geqslant \int_{E_0} |f| d\mu \geqslant \int_{E_0} n d\mu, \ \forall n \in \N.\]
    При этом
    \[\int_{E_0} n d\mu = n \cdot \int_{E_0} d\mu = n \cdot \mu(E_0).\]
    Тогда
    \[0 \leqslant \mu(E_0) \leqslant \frac{1}{n} \int_E |f| d\mu \to +0, n \to \infty.\]
\end{proof} 

\begin{statement}
    Пусть $E \in \mathfrak{M}$, пусть $\int_E |f| d\mu = 0$.
    Тогда $f = 0$ $\mu$-почти всюду на $E$.
\end{statement} 
\begin{proof}
    Предположим противное, то есть пусть существует $e \in E$ такое, что $\mu(e) > 0$ и $|f| > 0$ на $e$.
    Тогда по свойству для неотрицательных функций
    \[\int_e |f| d\mu > 0,\]
    чего не может быть.
\end{proof} 

\begin{statement}
    Пусть $f$ — ограниченная измеримая функция, $E \in \mathfrak{M}$ $\mu(E) < +\infty$.
    Тогда 
    \[\int_E f d\mu < +\infty.\]
\end{statement} 
\begin{proof}
    Пусть $\mu(E) < +\infty$.
    Поскольку $f$ ограничена и измерима, то существует константа $C > 0$ такая, что
    \[|f(x)| \leqslant C \ \forall x \in E.\]
    Тогда по свойствам интеграла Лебега для неотрицательных функций:
    \[\int_E |f| d\mu \leqslant \int_E C d\mu \leqslant C \cdot \mu(E) < +\infty \Longleftrightarrow f \text{интегрируема на } E.\]
\end{proof} 

\begin{statement}
    Пусть $f = g$ $\mu$-почти всюду на $E \in \mathfrak{M}$, где $f,g$ — измеримые функции.
    Тогда
    \begin{enumerate}
        \item $f$ интегрируема по Лебегу на $E$ тогда и только тогда, когда $g$ интегрируема по Лебегу на $E$.
        \item При этом, если обе интегрируемы, то
        \[\int_E f d\mu = \int_E g d\mu.\]
    \end{enumerate}
\end{statement} 
\begin{proof}
    $f$ интегрируема на $E$ тогда и только тогда, когда $|f|$ интегрируем на $E$.
    $g$ интегрируема на $E$ тогда и только тогда, когда $|g|$ интегрируем на $E$.
    Рассмотрим
    \[e := \{x \in E \ | \ f(x) \neq g(x)\}.\]
    Поскольку функции $\mu$-почти всюду совпадают, то $\mu(e) = 0$.
    Тогда по свойству интеграла для неотрицательных функций:
    \[\int_e |f| d\mu = \int_e |g| d\mu = 0,\]
    \[\int_{E \setminus e} |f| d\mu = \int_E |f| d\mu\]
    \[\int_{E \setminus e} |g| d\mu = \int_E |g| d\mu.\]
    — в силу аддитивности по множеству.
    Интегралы на $E \setminus e$ совпадают, поэтому
    \[\int_E |f| d\mu = \int_{E \setminus e} |f| d\mu = \int_{E \setminus e} |g| d\mu = \int_E |g| d\mu.\]
\end{proof} 

\begin{statement}
    Пусть $E, E' \in \mathfrak{M}$ такие, что $\mu(E \bigtriangleup E') = 0$.
    Тогда для любой измеримой на $X$ функции $f$ верно следующее:
    \begin{enumerate}
        \item $f$ интегрируема на $E$ тогда и только тогда, когда $f$ интегрируема на $E'$.
        \item Если $f$ интегрируема на $E$, то
        \[\int_E f d\mu = \int_{E'} f d\mu.\]
    \end{enumerate}
\end{statement} 
\begin{proof}
    Сведём к предыдущему случаю если
    \[f = \chi_E f, \tab g = \chi_{E'} f.\]
\end{proof} 

Напомним, что у нас изначально имеется тройка $(X, \mathfrak{M}, \mu)$ — абстрактное пространство с мерой.
Но мера не обязательно полна.
А если мера не полна, то могут быть подмножества множества меры ноль, которые неизмеримы.
Для таких случаев удобно ввести обобщённое понятие измеримости.

\begin{definition}
    Будем говорить, что $f$ \textit{измерима в широком смысле на $E$}, если существует $E' \subset E$ такое, что
    \[\mu(E \setminus E') = 0\]
    и $f$ измерима на $E'$.
\end{definition} 

Без ограничения общности считаем, что $f \equiv 0$ на $E \setminus E'$.

\begin{definition}
    \[\widetilde{L}_1 (E, \mu) := \left\{f: E \to \overline{\R} \ | \ f \text{ измерима в широком смысле на } E, \int_E |f| d\mu < +\infty\right\}.\]
\end{definition} 

\begin{statement}
    Пусть $f,g \in \widetilde{L}_1(E)$.
    Пусть $f \leqslant g$ $\mu$-почти всюду на $E$.
    Тогда
    \[\int_E f d\mu \leqslant \int_E g d\mu.\]
\end{statement} 
\begin{proof}
    Разложим $f$ и $g$:
    \[f = f_+ - f_-, \tab g = g_+ - g_-,\]
    тогда
    \[f_+ + g_- \leqslant g_+ + f_-\]
    $\mu$-почти всюду.

    Данное неравенство осмысленно, поскольку из интегрируемости вытекает конечность почти всюду на $E$.
    Тогда в силу монотонности интеграла
    \[\int_E (f_+ + g_-) d\mu \leqslant \int_E (g_+ + f_-) d\mu,\]
    а так как $f$ и $g$ интегрируемы по Лебегу на $E$, то
    \[\int_E f_+ d\mu - \int_E f_- d\mu \leqslant \int_E g_+ d\mu - \int_E g_- d\mu.\]
\end{proof} 

\begin{statement}[Аддитивность по функциям]
    Пусть $E \in \mathfrak{M}$, $f,g \in \widetilde{L}_1(E)$.
    Тогда
    \[\int_E (f+g) d\mu = \int_E f d\mu + \int_E g d\mu.\]
\end{statement} 
\begin{proof}
    \[f = f_+ - f_-, \tab g = g_+ - g_-.\]
    Пусть 
    \[h = f + g = (f_+ - f_-) + (g_+ - g_-) = (f_+ + g_+) - (f_- + g_-)\]
    — можем переставить слагаемые, поскольку из $f,g \in \widetilde{L}_1(E)$ следует, что они интегрируемы, значит, почти всюду конечны.
    С другой стороны,
    \[h = h_+ - h_-.\]
    Получается,
    \[h_+ + f_- + g_- = f_+ + g_+ + h_-,\]
    следовательно,
    \[\int_E (h_+ + f_- + g_-) d\mu = \int_E (f_+ + g_+ + h_-) d\mu,\]
    а в силу линейности для неотрицательных функций
    \[\int_E h_+ d\mu + \int_E f_- d\mu + \int_E g_- d\mu
    =
    \int_e h_- + \int_E f_+ d\mu + \int_E g_+ d\mu.\]
    В силу конечности всех слагаемых и определения интеграла Лебега получаем утверждение.
\end{proof} 

\begin{statement}[Однородность интеграла Лебега]
    Пусть $E \in \mathfrak{M}$, $f \in \widetilde{L}_1(E,\mu)$.
    Тогда для любого $\alpha \in \R$ верно
    \[\int_E \alpha \cdot f d\mu = \alpha \cdot \int_E f d\mu.\]
\end{statement} 
\begin{proof}
    Пусть $\alpha \geqslant 0$.
    Тогда для $f = f_+ - f_-$ следует $$\alpha \cdot f = \alpha \cdot f_+ - \alpha \cdot f_-.$$
    Из положительной однородности интеграла Лебега для неотрицательных функций:
    \[\alpha \cdot \int_E f_+ d\mu = \int_E \alpha \cdot f_+ d\mu,\]
    \[\alpha \cdot \int_E f_- d\mu = \int_E \alpha \cdot f_- d\mu,\]
    тогда
    \[\int_E \alpha \cdot f d\mu
    = \alpha \cdot \int_E f_+ d\mu - \alpha \cdot \int_E f_- d\mu
    = \alpha \cdot \int_E f d\mu.\]
    Если же $\alpha < 0$, то достаточно рассмотреть случай $\alpha = -1$, поскольку любое отрицательное число $\beta = (-1) \cdot \alpha$.
    В этом случае
    \[(-f)_+ = f_-, \tab (-f)_- = f_+,\]
    Из предыдущих рассуждений:
    \begin{multline*}
    \int_E -f d\mu = \int_E (-f)_+ d\mu - \int_E (-f)_- d\mu
    =\\= - \int_E f_+ d\mu + \int_E f_- d\mu
    =\\= - \left(\int_E f_+ d\mu - \int_E f_- d\mu\right)
    = - \int_E f d\mu.
    \end{multline*} 
\end{proof} 

\begin{statement}
    Если $f_1, \dots, f_N \in \widetilde{L}_1(E, \mu)$ и $\alpha_1, \dots, \alpha_N \in \R$.
    Тогда
    \[\sum\limits_{i=1}^{N} \alpha_i f_i \in \widetilde{L}_1(E, \mu).\]
    Кроме того,
    \[\int_E \left(\sum\limits_{i=1}^{N} \alpha_i f_i\right) d\mu = \sum\limits_{i=1}^{N} \left(\alpha_i \cdot \int_E f_i d\mu\right).\]
\end{statement} 
\begin{proof}
    Доказательство вытекает из предыдущих свойств.
\end{proof} 